% ============================================================================
% instantiations/_TEMPLATE/set.tex — INSTITUTION SET (Layer 2)
% ----------------------------------------------------------------------------
% This file is the single localization entry point for ONE adopting
% institution. It is the Layer-2 file of the suite localization model:
%
%     localization.sty                         (\setloc / \loc machinery)
%       Layer 0: shared/localization-defaults.tex   (generic role phrases)
%         Layer 1: categories/<ARCHETYPE>.tex        (this institution's class)
%           Layer 2: THIS FILE                        (concrete institution values)
%             document body
%
% A render wrapper (render/*.tex) loads Layer 0, then \input's THIS file. This
% file in turn \input's the chosen category archetype (Layer 1), then applies
% the institution's own concrete \setloc values and any local deltas (Layer 2).
% Each later \setloc overrides any key an earlier layer set; keys left untouched
% keep the value from below.
%
% HOW TO FILL THIS FILE (see ./README.md for the full walk-through):
%   1. Pick your archetype with categories/README.md (the §2 decision tree,
%      run after the SUITSOLUTION Chapter 4 self-assessment).
%   2. Set \instcategory below to the chosen archetype file name (no path,
%      no .tex), e.g. EU-PUB-RI, US-PUB-R1, US-PRIV-R1, UK-RES, LATAM-PUB,
%      AFR-PUB.
%   3. Fill the country-specific NAME keys your archetype leaves at default.
%   4. Add any LOCAL DELTAS the decision tree told you to record (a key your
%      archetype binds generically but that differs for your institution).
%   5. Build the four documents from render/ and deposit (see README.md §4).
%
% Keys are kebab-case. The full key catalogue is shared/localization-defaults.tex
% (Layer 0) and categories/README.md §3 (the per-archetype comparison matrix).
% Calling an undefined key fails the build loudly, so a forgotten value is
% always caught at compile time.
% ============================================================================

% --- 1. CHOSEN CATEGORY ARCHETYPE (Layer 1) ---------------------------------
% Set \instcategory to your archetype's file name (without path or extension):
% EU-PUB-RI, US-PUB-R1, US-PRIV-R1, UK-RES, LATAM-PUB or AFR-PUB.
%
% The render wrapper has already defined \catdir as the relative path from its
% own location to the repository's categories/ directory, so the \input below
% resolves to categories/<ARCHETYPE>.tex from whichever document is building.
% (Do NOT redefine \catdir here.)
\newcommand{\instcategory}{EU-PUB-RI}% <-- REPLACE with your archetype
\input{\catdir/\instcategory}% -> categories/<ARCHETYPE>.tex

% ============================================================================
% 2. COUNTRY-SPECIFIC NAME KEYS (left at default by every archetype)
% ----------------------------------------------------------------------------
% Your archetype binds the STRUCTURAL and REGULATORY keys generically but
% deliberately leaves these per-country NAME keys for you to fill. Replace each
% placeholder with your institution's concrete value. (UK-RES fixes the NREN to
% Janet and US-PRIV-R1 fixes it to Internet2 — leave nren-name alone there.)
% ============================================================================
\setloc{nren-name}{REPLACE: your national research-and-education network, by name}
\setloc{national-dpa}{REPLACE: your national data-protection authority, by name}
\setloc{national-gov-csirt}{REPLACE: your national / governmental CSIRT, by name}
% US archetypes only — the regional R&E network security team:
% \setloc{nren-csirt}{REPLACE: your state/regional R&E network security team}

% ============================================================================
% 3. EDITION IDENTITY (banner REQUIRED) — INSTITUTION IDENTITY (optional)
% ----------------------------------------------------------------------------
% edition-name is REQUIRED for every deposit: it is rendered prominently on
% each title page (\suiteditionline) and in \suitmetablock, so every PDF says
% at a glance which edition it belongs to. An anonymous edition names itself
% by its archetype instead of the institution.
% ============================================================================
\setloc{edition-name}{REPLACE: Edition of the University of Example (or: Edition for a continental-European public research university)}

% Uncomment and fill to render a named edition rather than a generic-by-class
% one. Leaving these commented keeps the institution anonymous behind its
% archetype, which is a legitimate deposit.
% \setloc{institution-name}{REPLACE: the University of Example}
% \setloc{institution-type}{REPLACE: a public, research-intensive university established under the law of ...}
% \setloc{executive-leadership}{REPLACE: the rectorate / the President and Provost / the Vice-Chancellorate}
% \setloc{governing-board}{REPLACE: the Board of Governors / Council / Board of Regents}
% \setloc{academic-governing-bodies}{REPLACE: the Senate / Academic Council}

% ============================================================================
% 4. LOCAL DELTAS (record every deviation the decision tree surfaced)
% ----------------------------------------------------------------------------
% If your institution does not match the chosen archetype exactly (an NREN-tier
% delta, a governance delta, a derived regulatory row from the mapping method
% in shared/regulatory-applicability-EN.tex), re-bind the affected structural
% or regulatory key HERE so it overrides the archetype. Document WHY in a
% comment — the delta is a real finding, not noise.
% ============================================================================
% Example NREN-tier delta (EU institution reaching GÉANT via a regional aggregator):
% \setloc{nren-name}{the regional R&E aggregator <name>, reaching GÉANT through the national NREN <name>}
%
% Example derived regulatory row (jurisdiction not among the worked regimes):
% \setloc{data-protection-regime}{the national data-protection law <citation> (derived via the sec:reg-mapping Step 1–7 method)}

\endinput
